\chapter{形状、概率路径与监督矩}
\label{app:shape}
\section{覆盖界与目标半径}
\begin{theorem}[正形状的覆盖与刷新后信号范围]
固定正分布 $q,T$，$0<\kappa\le1,0<\eta<1$。幂均值的归一化常数和目标满足
\begin{equation}
(1-\eta)^{1/\kappa}+\eta^{1/\kappa}\le Z\le1,\qquad
m_a\ge(1-\eta)^{1/\kappa}q_a+\eta^{1/\kappa}T_a.
\end{equation}
因此 $\KL(q\Vert m)\le-\log(1-\eta)/\kappa$、$\KL(T\Vert m)\le-\log\eta/\kappa$。精确刷新 $q^+=m$ 后，
$-c_{\kappa,\eta}\le A(T_a/q_a^+)\le\log(2-\eta)/(\kappa\eta)$。
\end{theorem}
\begin{proof}
对 $0<\kappa\le1$，幂均值不等式给 $h_a\le(1-\eta)q_a+\eta T_a$，所以 $Z\le1$。令 $r=1/\kappa\ge1$，非负 $x,y$ 满足 $(x+y)^r\ge x^r+y^r$；例如对固定 $x$，两侧之差对 $y$ 的导数非负，且在 $y=0$ 为零。因此
\begin{equation}
(1-\eta)^{1/\kappa}+\eta^{1/\kappa}\le Z\le1,\qquad
m_a\ge(1-\eta)^{1/\kappa}q_a+\eta^{1/\kappa}T_a.
\end{equation}
逐坐标取对数并按相应分布平均，得到
\begin{equation}
\KL(q\Vert m)\le-\kappa^{-1}\log(1-\eta),\quad
\KL(T\Vert m)\le-\kappa^{-1}\log\eta.
\end{equation}
固定教师的精确刷新 $q^+=m$ 还给后续概率比 $(T_a/q_a^+)^\kappa\le1/\eta$，所以下一轮
\begin{equation}
-c_{\kappa,\eta}\le A_{\kappa,\eta}(T_a/q_a^+)
\le\frac{\log(2-\eta)}{\kappa\eta}.
\end{equation}
这个双侧界依赖精确刷新产生的覆盖；任意初始锚点只有单侧下界。
\end{proof}

\begin{proposition}[目标到锚点的半径]
固定 $q,T,\kappa$。对任意凸 $\varphi$、$\varphi(1)=0$，$D_\varphi(m_{\kappa,\eta}\Vert q)$ 随 $\eta\in[0,1]$ 非减。若 $q\ne T$ 且 $\varphi$ 严格凸，内部严格递增。
\end{proposition}
\begin{proof}
取 $0<\eta_1<\eta_2<1$，令 $\rho=T/q$、$u_i=m_{\eta_i}/q$。$u_1$ 随 $\rho$ 递增，而
\begin{equation}
\frac{m_{\eta_2}}{m_{\eta_1}}
=\frac{Z_{\eta_1}}{Z_{\eta_2}}
\left(\frac{1-\eta_2+\eta_2\rho^\kappa}
{1-\eta_1+\eta_1\rho^\kappa}\right)^{1/\kappa}
\end{equation}
也严格递增：括号对 $\rho^\kappa$ 的导数为 $(\eta_2-\eta_1)/(1-\eta_1+\eta_1\rho^\kappa)^2>0$。几何端点的比值为常数乘 $\rho^{\eta_2-\eta_1}$。该比值在 $m_{\eta_1}$ 下均值一，故存在交点 $\rho_\star$ 使比值等于一。

记 $u_\star=u_1(\rho_\star)$。低于交点时 $u_2\le u_1\le u_\star$，高于交点时 $u_2\ge u_1\ge u_\star$。取 $b\in\partial\varphi(u_\star)$，则 $\varphi(u)-bu$ 在 $u_\star$ 两侧向外非减，逐点得到 $\varphi(u_2)-bu_2\ge\varphi(u_1)-bu_1$。按 $q$ 求和，因两组 $u$ 的均值均为一，线性项抵消。严格凸与不同目标保证至少一项严格；端点由连续性得到。
\end{proof}

\section{精确目标同时改善凸生成散度}
\begin{theorem}[共同交点的散度进展]
对任意凸 $\varphi$、$\varphi(1)=0$，精确幂均值目标满足
\begin{equation}
D_\varphi(m\Vert T)\le D_\varphi(q\Vert T).
\end{equation}
\end{theorem}
\begin{proof}
写 $s_a=q_a/T_a$、$m_a=T_aG(s_a)/Z$，其中 $G(s)=[(1-\eta)s^\kappa+\eta]^{1/\kappa}$，几何端点为 $G(s)=s^{1-\eta}$。$G$ 递增，$G(s)/s$ 递减，因为
\begin{equation}
\frac d{d\log s}\log[G(s)/s]
=-\frac{\eta s^{-\kappa}}{1-\eta+\eta s^{-\kappa}}<0.
\end{equation}
又有 $Z=\sum_aq_aG(s_a)/s_a$，所以存在 $s_\star$ 使 $G(s_\star)/s_\star=Z$。若 $s_a<s_\star$，则 $s_a\le G(s_a)/Z\le s_\star$；若 $s_a>s_\star$，不等号反向排列。取 $\varphi$ 在 $s_\star$ 的次梯度 $b$，函数 $\varphi(s)-bs$ 向 $s_\star$ 移动时非增。逐项比较后按 $T$ 求和，归一化消去线性项，得到结论。
\end{proof}

\begin{proposition}[精确更新的可逆 Markov 表示]
对固定正分布 $q,T$ 和精确幂均值目标 $m$，存在行随机矩阵 $P$，满足 $TP=T$、$qP=m$，且关于 $T$ 详细平衡：$T_aP_{ab}=T_bP_{ba}$。
\end{proposition}
\begin{proof}
这个更新还可由保持 $T$ 的可逆 Markov 核表示。令 $\delta_a=m_a-q_a$，正质量接收集合为 $L$，输出集合为 $H$，总转移量 $D=\sum_L\delta_a$。若 $D=0$ 取恒等核；否则对 $a\in L,b\in H$，定义
\begin{equation}
w_{ab}=w_{ba}=\frac{\delta_a(-\delta_b)}{D(s_b-s_a)},\qquad
P_{ab}=w_{ab}/T_a,
\end{equation}
其余非对角项为零，对角项补足行和。由共同交点，$\delta_a/T_a\le s_\star-s_a$、$(-\delta_b)/T_b\le s_b-s_\star$，故 $\sum_bw_{ab}\le T_a$、$\sum_aw_{ab}\le T_b$，核非负。详细平衡 $T_aP_{ab}=T_bP_{ba}$ 保证 $TP=T$；而接收坐标有 $(qP)_a-q_a=\sum_bw_{ab}(s_b-s_a)=\delta_a$，输出坐标同理，故 $qP=m$。这是数学表示，无需实际构造词表平方大小的矩阵。
\end{proof}

\section{教师中心的性质}
\begin{proposition}[教师中心的锐统一界]
对 $0<\kappa\le1,0<\eta<1$，任意正分布满足
\begin{equation}
0\le K_T(q,T)\le c_{\kappa,\eta}=-\frac{\log(1-\eta)}{\kappa\eta}.
\end{equation}
至少二动作时，该上界可由严格正分布逼近。
\end{proposition}
\begin{proof}
写 $\rho=T/q$，则 $A(1)=0$、$A'>0$ 且 $A''\le0$。Jensen 给出 $\E_qA(\rho)\le A(\E_q\rho)=0$；结合 $A\ge-c$，
\begin{equation}
0\le K_T=-\E_qA\le c_{\kappa,\eta}.
\end{equation}
上界为锐上确界：取 $q=(\epsilon,1-\epsilon)$、$T=(1-\epsilon^2,\epsilon^2)$。常见动作的 $-A\to c$，稀有动作的贡献为 $O(\epsilon|\log\epsilon|)\to0$。
\end{proof}

\begin{proposition}[中心对两个参数的单调性]
固定 $q,T>0$，教师中心 $K_T$ 随 $\kappa$ 非增、随 $\eta$ 非减，参数范围为 $0<\kappa\le1,0<\eta<1$。该结论不对监督方差给出同样排序。
\end{proposition}
\begin{proof}
固定概率比时，$A$ 随 $\kappa$ 非减：函数 $L(u)=\log(1-\eta+\eta e^{u\log\rho})$ 凸且 $L(0)=0$，故 $L(u)/u$ 在 $u>0$ 非减。另一方面，令 $y=1+\eta(\rho^\kappa-1)$，有
\begin{equation}
\partial_\eta A=\frac{1-y^{-1}-\log y}{\kappa\eta^2}\le0.
\end{equation}
不等式来自 $\log y\ge1-y^{-1}$。因此 $K_T$ 随 $\kappa$ 非增，随 $\eta$ 非减。标量方差没有由此得到同一排序。
\end{proof}

\begin{proposition}[中心的端点与局部尺度]
固定正分布时，$K_T\to(1-\kappa)D_\kappa(T\Vert q)$ 当 $\eta\downarrow0$，而 $\eta\uparrow1$ 时 $K_T\to\KL(q\Vert T)$。若全部 $\rho=T/q$ 一致趋近一，则
\begin{equation}
K_T=\tfrac12[1-\kappa(1-\eta)]\E_q(\rho-1)^2+O(\E_q|\rho-1|^3).
\end{equation}
\end{proposition}
\begin{proof}
固定概率比时，$\log[1+\eta(\rho^\kappa-1)]/\eta\to\rho^\kappa-1$。有限求和可逐项取极限，所以
\begin{align}
\lim_{\eta\downarrow0}K_T
&=\frac{1-\E_q\rho^\kappa}{\kappa}=(1-\kappa)D_\kappa(T\Vert q),\\
K_T&=\frac{1-\kappa(1-\eta)}2\E_q(\rho-1)^2
+O(\E_q|\rho-1|^3).
\end{align}
第一式对 $\kappa=1$ 取零。为证明第二式，直接计算
\begin{equation}
A(1)=0,\qquad A'(1)=1,\qquad A''(1)=\kappa(1-\eta)-1=-c_0.
\end{equation}
在一附近 Taylor 展开为 $A(1+\delta)=\delta-c_0\delta^2/2+O(|\delta|^3)$。按 $q$ 取负期望，一次项由 $\E_q(\rho-1)=0$ 消失，得到局部尺度。$\eta\uparrow1$ 时逐项 $A(\rho)\to\log\rho$，所以 $K_T\to\KL(q\Vert T)$；统一有界性不延续到该端点。
\end{proof}

\section{固定强度的矩界与小步阈值}
\begin{proposition}[固定强度的维数无关二阶矩界]
令 $C_\star=\sup_{u>0}\log^2(1+u)/u<\infty$。对 $0<\kappa\le1,0<\eta<1$，未缩放信号 $f=\eta A$ 满足
\begin{equation}
\E_qf^2\le[\log^2(1-\eta)+C_\star\eta]/\kappa^2.
\end{equation}
标准化信号 $A$ 的相应上界为右侧除以 $\eta^2$。
\end{proposition}
\begin{proof}
定义 $C_\star=\sup_{u>0}\log^2(1+u)/u<\infty$。函数在零与无穷远都趋零，故上确界有限。对 $f=\eta A$，在 $\rho<1$ 上 $|f|\le|\log(1-\eta)|/\kappa$；在 $\rho\ge1$ 上
\begin{equation}
0\le f\le\kappa^{-1}\log(1+\eta\rho^\kappa),\qquad
f^2\le C_\star\eta\rho^\kappa/\kappa^2.
\end{equation}
凹性给出 $\E_q\rho^\kappa\le1$，所以
\begin{equation}
\E_qf^2\le\frac{\log^2(1-\eta)+C_\star\eta}{\kappa^2}.
\label{eq:raw-moment}
\end{equation}
标准化监督 $A$ 的界需再除以 $\eta^2$。
\end{proof}

\begin{theorem}[小步矩阈值]
固定 $r\ge1$、$0<\kappa\le1$。在所有有限正概率分布上取上确界，有
\begin{equation}
\sup_{q,T}\E_q|A|^r
=\Theta_{\kappa,r}(\eta^{\min\{r,1/\kappa\}-r}),\qquad
\sup_{q,T}\E_q|C_qA|^r
=\Theta_{\kappa,r}(\eta^{\min\{r,1/\kappa\}-r})
\end{equation}
当 $\eta\downarrow0$。$\kappa=0$ 对任何正 $\eta$ 都没有有限统一矩界。
\end{theorem}
\begin{proof}
先估计 $f$。负区域为 $O_{\kappa,r}(\eta^r)$。若 $r\kappa\le1$，用 $\log(1+x)\le x$ 和 $\E_q\rho^{r\kappa}\le1$ 得正区域也是 $O(\eta^r)$。若 $r\kappa>1$，令 $s=1/\kappa<r$。函数 $\log^r(1+x)/x^s$ 在零处按 $x^{r-s}$ 趋零，在无穷处也趋零，故有有限上界 $C_{r,s}$。于是
\begin{equation}
\E_q[|f|^r\mathbf1_{\rho\ge1}]
\le\kappa^{-r}C_{r,s}\eta^{1/\kappa}\E_q\rho
=O(\eta^{1/\kappa}).
\end{equation}
除以 $\eta^r$ 得未中心化上界。由 $|X-\E X|^r\le2^{r-1}(|X|^r+|\E X|^r)$ 及 Jensen，中心化矩不超过 $2^r\E|X|^r$。

当 $r\kappa\le1$，任意固定不同的二元正分布使 $A\to(\rho^\kappa-1)/\kappa$，其原始与中心化矩均为严格正常数，给出下界。当 $r\kappa>1$，取 $q_1=\eta^{1/\kappa}/2$、$T_1=1/2$。稀有动作有 $A_1=\log(2-\eta)/(\kappa\eta)$，另一动作 $A_2=O(1)$。所以 $q_1|A_1|^r=\Omega(\eta^{1/\kappa-r})$；中心化后的稀有值为 $(1-q_1)(A_1-A_2)$，具有同阶贡献。几何端点取固定非退化 $q$，令一个教师概率趋零，即使中心化，正质量动作上的对数差也无界。
\end{proof}

二阶矩阈值是 $\kappa\le1/2$，四阶矩阈值是 $\kappa\le1/4$。固定分布的小步 Taylor 展开不能替代这类允许分布随 $\eta$ 改变的最坏情形论证。

\section{误差相关方差与锐常数}
\begin{theorem}[误差相关的监督方差上界]
若 $0<\kappa\le1/2,0<\eta<1$，则
\begin{align}
V_{\kappa,\eta}&\le\frac{1-a_\kappa^2}{\kappa^2(1-\eta)^2},\qquad a_\kappa=\E_q(T/q)^\kappa,\\
V_{\kappa,\eta}&\le\frac{D_\kappa(T\Vert q)}{\kappa(1-\eta)^2}
\le\frac{K_T}{\kappa(1-\kappa)(1-\eta)^2}.
\end{align}
\end{theorem}
\begin{proof}
令 $a_\kappa=\E_q\rho^\kappa$。当 $0<\kappa\le1/2$，函数 $h(x)=\log(1-\eta+\eta x)/(\kappa\eta)$ 的 Lipschitz 常数不超过 $1/[\kappa(1-\eta)]$。独立同分布 $X,X'$ 满足 $\Var h(X)=\frac12\E[h(X)-h(X')]^2$，故
\begin{equation}
V_{\kappa,\eta}\le\frac{\Var_q(\rho^\kappa)}{\kappa^2(1-\eta)^2}
\le\frac{1-a_\kappa^2}{\kappa^2(1-\eta)^2}.
\end{equation}
第二步用 $\E_q\rho^{2\kappa}\le1$。

为得到与教师误差相关的界，令 $a(u)=\E_q\rho^u$。其二阶导数为 $\E_q[\rho^u\log^2\rho]\ge0$，所以对 $2\kappa\in[\kappa,1]$，
\begin{equation}
a(2\kappa)\le\frac{1-2\kappa}{1-\kappa}a(\kappa)+\frac{\kappa}{1-\kappa}.
\end{equation}
减去 $a(\kappa)^2$ 并因式分解：
\begin{equation}
\Var_q(\rho^\kappa)
\le\frac{(1-a_\kappa)[\kappa+(1-\kappa)a_\kappa]}{1-\kappa}
\le\frac{1-a_\kappa}{1-\kappa}.
\end{equation}
代入散度定义得到式~\eqref{eq:error-variance-main}。最后一步利用 $\log(1+\eta x)\le\eta x$ 得 $K_T\ge(1-a_\kappa)/\kappa$。
\end{proof}

\begin{theorem}[相对教师误差的锐小步方差常数]
设 $A_\kappa^0=(\rho^\kappa-1)/\kappa$、$K_T^0=-\E_qA_\kappa^0$。在至少二动作的正分布类中，对 $0<\kappa\le1/2$，
\begin{equation}
\sup_{q\ne T}\frac{\Var_qA_\kappa^0}{D_\kappa(T\Vert q)}=\frac1\kappa,\qquad
\sup_{q\ne T}\frac{\Var_qA_\kappa^0}{K_T^0}=\frac1{\kappa(1-\kappa)}.
\end{equation}
其余形状没有有限统一常数，故两种准则均由 $\kappa=1/2$ 唯一最小化。
\end{theorem}
\begin{proof}
在小步极限 $A^0_\kappa=(\rho^\kappa-1)/\kappa$，
\begin{equation}
\sup_{q\ne T}\frac{\Var_qA^0_\kappa}{D_\kappa(T\Vert q)}=\frac1\kappa,\qquad
\sup_{q\ne T}\frac{\Var_qA^0_\kappa}{K_T^0}=\frac1{\kappa(1-\kappa)}.
\end{equation}
上界已证。锐性取 $q=(v,1-v)$、$T=(v^2,1-v^2)$，此时 $\rho=(v,1+v)$。当 $v\downarrow0$，
\begin{equation}
a_\kappa=1-(1-\kappa)v+o(v),\quad
\Var_qA^0_\kappa=v/\kappa^2+o(v),\quad
D_\kappa=v/\kappa+o(v).
\end{equation}
两比值逼近所列常数。对 $\kappa>1/2$，取 $q=(v,1-v),T=(1/2,1/2)$，方差按 $v^{1-2\kappa}$ 发散；内部幂散度有界，$\kappa=1$ 的 KL 只按 $\log(1/v)$ 增长。$\kappa=0$ 用固定非对称 $q$ 与趋零教师概率，方差按对数平方增长，而 KL 按一次对数增长。故相对准则的唯一最优点为 $\kappa=1/2$。
\end{proof}

\section{绝对最坏方差的不同最优点}
\begin{theorem}[绝对最坏小步方差及其最优形状]
对 $0<\kappa<1/2$，令 $b=(1-2\kappa)/[2(1-\kappa)]$，则
\begin{equation}
\sup_{q,T}\Var_q A_\kappa^0
=\frac{b^{1-2\kappa}(1-b)}{\kappa^2}.
\end{equation}
在 $\kappa=1/2$，上确界为四。该绝对方差准则的唯一最优形状满足
$\log[(1-2\kappa)/(2(1-\kappa))]=-1/\kappa$，数值约为 $0.441928$。
\end{theorem}
\begin{proof}
对 $0<\kappa<1/2$，令 $X=\rho^\kappa$、$r=1/\kappa>2$，则 $\E_qX^r=1$。Hölder 给出
\begin{equation}
\E X^2\le(\E X)^{(r-2)/(r-1)}.
\end{equation}
具体地，将 $X^2=X^{(r-2)/(r-1)}(X^r)^{1/(r-1)}$，取共轭指数 $(r-1)/(r-2)$ 与 $r-1$ 即得。令 $a=\E X$、$\gamma=(1-2\kappa)/(1-\kappa)$，则 $\Var X\le a^\gamma-a^2$。最大值处 $a^{2-\gamma}=\gamma/2$，得到
\begin{equation}
\sup_{q,T}\Var A^0_\kappa
=C_{\rm vel}(\kappa)
=\frac{b^{1-2\kappa}(1-b)}{\kappa^2},
\qquad b=\frac{1-2\kappa}{2(1-\kappa)}.
\end{equation}
取 $q=(b,1-b)$、$T\to(1,0)$，Hölder 与驻点条件同时取等，证明锐性。$\kappa=1/2$ 的上确界为四。对最优 $b$ 使用包络求导，
\begin{equation}
\frac d{d\kappa}\log C_{\rm vel}=-2/\kappa-2\log b.
\end{equation}
唯一根满足 $\log[(1-2\kappa)/(2(1-\kappa))]=-1/\kappa$。左侧加 $1/\kappa$ 严格递减，端点分别趋于正、负无穷，故根唯一，数值为 $\kappa_\star\simeq0.441928$。这个准则最小化绝对最坏方差，不等于相对误差准则，也不等于任务表现。
\end{proof}

\section{平方根流、局部速率与噪声可求和性}
\begin{proposition}[平方根自然梯度流的显式亲和度]
固定正教师，沿 $\kappa=1/2$ 的自然梯度流，令 $a=\sum\sqrt{qT}$。则 $\|v_{1/2}\|_q^2=4(1-a^2)$、$\dot a=1-a^2$。若初值不同于教师，
$a(t)=\tanh[t+\operatorname{artanh}a(0)]$，且 $D(0)e^{-2t}\le D_{1/2}(T\Vert q(t))\le D(0)e^{-t}$。
\end{proposition}
\begin{proof}
平方根点令 $a=\sum\sqrt{qT}$。小步 Fisher 速度满足 $\|v_{1/2}\|_q^2=4(1-a^2)$，且 $D_{1/2}=4(1-a)$。因此沿连续自然梯度流，
\begin{equation}
\dot a=1-a^2,\qquad a(t)=\tanh[t+\operatorname{artanh}a(0)]
\end{equation}
对 $q\ne T$ 成立；在 $q=T$ 处流恒定。又有 $\dot D=-(1+a)D$，故 $D(0)e^{-2t}\le D(t)\le D(0)e^{-t}$。离散精确更新的 Hellinger 收缩由正文式~\eqref{eq:hellinger-contraction-main} 证明。
\end{proof}

\begin{proposition}[幂均值的二阶展开与共同局部率]
固定正教师、$\kappa\in[0,1],0<\eta<1$。令 $q=T+\epsilon h$、$\sum h=0$、$S_h=\sum h_a^2/T_a$，则
\begin{equation}
m_a=T_a+(1-\eta)\epsilon h_a+
\frac{(\kappa-1)(1-\eta)\eta}{2}\epsilon^2
\left(h_a^2/T_a-T_aS_h\right)+O(\epsilon^3).
\end{equation}
对固定非零 $h$ 及一处具有正二阶曲率的三阶光滑 $f$-散度，更新前后教师散度比为 $(1-\eta)^2+O(\epsilon)$。
\end{proposition}
\begin{proof}
固定 $T>0$，写 $q=T+\epsilon h$、$\sum h=0$、$S_h=\sum h_a^2/T_a$。两次二项式展开并除以归一化常数，得到
\begin{equation}
m_a=T_a+(1-\eta)\epsilon h_a
+\frac{(\kappa-1)(1-\eta)\eta}{2}\epsilon^2
\left(\frac{h_a^2}{T_a}-T_aS_h\right)+O(\epsilon^3).
\end{equation}
未归一化展开中二阶项为 $(\kappa-1)(1-\eta)\eta\epsilon^2h_a^2/(2T_a)$；其总和正是需要减去的归一化项。几何端点直接展开 $q^{1-\eta}T^\eta$ 给出相同结果。任何在一处二阶曲率为正的光滑 $f$-散度都满足局部比值 $(1-\eta)^2+O(\epsilon)$。
\end{proof}

\begin{corollary}[精确外层的监督噪声可求和性]
固定教师、$0<\kappa\le1/2,0<\eta<1$，每轮精确刷新。则
\begin{equation}
\sum_{n=1}^\infty V_n\le
\frac{\KL(T\Vert q_0)}{\kappa(1-\kappa)(1-\eta)^2\eta}.
\end{equation}
平方根点还满足 $V_n\le4(1-\eta)^nH^2(q_0,T)/(1-\eta)^2$。该式控制上界包络，不断言实际方差逐步单调。
\end{corollary}
\begin{proof}
最后，固定 $0<\kappa\le1/2,\eta<1$ 并精确刷新，式~\eqref{eq:center-sandwich} 与三点求和给出
\begin{equation}
\sum_{n=1}^\infty V_n
\le\frac{\KL(T\Vert q_0)}{\kappa(1-\kappa)(1-\eta)^2\eta}.
\end{equation}
平方根点可用更直接的 Hellinger 界：
\begin{equation}
V_n\le\frac{4H^2(q_n,T)}{(1-\eta)^2}
\le\frac{4(1-\eta)^nH^2(q_0,T)}{(1-\eta)^2}.
\end{equation}
上界包络下降，不保证实际方差逐步单调；所有这些结论固定同一前缀、教师和参数。
\end{proof}

\section{固定 Fisher 速度下的比较}
\begin{proposition}[共同损失下的最陡形状]
\label{prop:fisher-shape}
固定有限正分布 $q\ne T$，令 $\rho=T/q$、$r=\log\rho$，
$w_\kappa=C_q(\rho^\kappa)/\kappa$，并以连续延拓定义 $w_0=C_qr$。
向量 $v_\kappa=q w_\kappa$ 是小步目标的切向量。将其按 Fisher 范数
$\|v\|_q^2=\sum_a v_a^2/q_a$ 归一化为 $u_\kappa$，则
\begin{align}
-D\KL(q\Vert T)[u_\kappa]
 &=\frac{\E_q[rw_\kappa]}{\sqrt{\E_qw_\kappa^2}}
 \le\sqrt{\Var_qr},\\
-D\KL(T\Vert q)[u_\kappa]
 &=\frac{\E_q[\rho w_\kappa]}{\sqrt{\E_qw_\kappa^2}}
 \le\sqrt{\Var_q\rho}.
\end{align}
第一项由 $\kappa=0$ 达到，第二项由 $\kappa=1$ 达到；不声称达到者唯一。
\end{proposition}
\begin{proof}
切向量满足 $\sum_a v_a=0$。逐坐标微分两个 KL，
其负方向导数分别为 $\sum_a r_av_a$ 与 $\sum_a\rho_av_a$。
代入 $u_\kappa=q w_\kappa/\sqrt{\E_qw_\kappa^2}$ 得等式。
由于 $\E_qw_\kappa=0$，分子可分别写成
$\E_q[(C_qr)w_\kappa]$ 和 $\E_q[(C_q\rho)w_\kappa]$。
Cauchy--Schwarz 不等式给上界；在两个所述端点，相应中心化向量正比例，故取等。
这里统一的是 Fisher 速度和被比较的损失，与统一原始步长、最小化监督方差是不同的比较条件。
\end{proof}

\section{平方根目标与平方根拟合的联合界}
\begin{proposition}[初始拟合系数的统一二阶矩]
\label{prop:root-joint}
取 $\kappa=1/2$ 的目标 $m$，$0<\eta\le1$，并使用平方根散度拟合。
在初始点 $p=q$，令 $a=\sum_i\sqrt{q_iT_i}$、
$b=\sum_i\sqrt{q_im_i}$。若归一化拟合系数为
$G_i=2(b-\sqrt{m_i/q_i})/\eta$，则
\begin{equation}
\E_qG^2=\frac{4(1-a^2)}{Z}
\le\frac4{(1-\eta)^2+\eta^2}\le8,\qquad
Z=(1-\eta)^2+\eta^2+2\eta(1-\eta)a.
\end{equation}
若 score 范数至多为 $S$，相应单动作梯度的二阶矩至多为 $8S^2$。
\end{proposition}
\begin{proof}
由 $\sqrt{m_i}=[(1-\eta)\sqrt{q_i}+\eta\sqrt{T_i}]/\sqrt Z$，
有 $b=(1-\eta+\eta a)/\sqrt Z$。
又 $\E_q(m/q)=1$，所以 $\E_qG^2=4(1-b^2)/\eta^2$。
直接相减得到 $Z-(1-\eta+\eta a)^2=\eta^2(1-a^2)$。
因 $0<a\le1$，丢弃 $Z$ 的非负交叉项并用
$(1-\eta)^2+\eta^2\ge1/2$ 得界。最后逐点使用
$\|G_i\psi_i\|^2\le S^2G_i^2$。
此处拟合形状也是平方根；该系数不是 BPO 的对数残差。
\end{proof}

\begin{proposition}[有限步最坏方差的小步极限]
\label{prop:finite-variance-limit}
设 $0<\kappa\le1/2$，$V_{\kappa,\eta}(q,T)=\Var_q A_{\kappa,\eta}(T/q)$。
以前文的锐常数
$C_{\rm vel}(\kappa)=\sup_{q,T}\Var_q(\rho^\kappa)/\kappa^2$ 记最坏速度方差，则
\begin{equation}
\sup_{q,T}V_{\kappa,\eta}(q,T)
\le\frac{C_{\rm vel}(\kappa)}{(1-\eta)^2},\qquad
\lim_{\eta\downarrow0}\sup_{q,T}V_{\kappa,\eta}(q,T)
=C_{\rm vel}(\kappa).
\end{equation}
固定 $q,T$ 时，平方根端点的极限为
$4(1-a^2)=H^2(q,T)[4-H^2(q,T)]$，其中 $H^2=2-2a$。
\end{proposition}
\begin{proof}
将 $A$ 看成 $x=\rho^\kappa$ 的函数，其导数为
$1/[\kappa(1-\eta+\eta x)]\le1/[\kappa(1-\eta)]$。
对独立同分布的 $X,X'$，恒等式
$\Var h(X)=\E(h(X)-h(X'))^2/2$ 将 Lipschitz 界转为方差界。
固定正分布时逐点 Taylor 展开给
$A_{\kappa,\eta}(\rho)\to(\rho^\kappa-1)/\kappa$。
对任意 $\epsilon>0$，先选择一个速度方差超过
$C_{\rm vel}-\epsilon$ 的固定分布对，再令 $\eta\downarrow0$，
得到下极限至少为 $C_{\rm vel}-\epsilon$；上界给上极限不超过
$C_{\rm vel}$。最后令 $\epsilon\downarrow0$。
平方根公式由 $\E_q\rho=1$ 和 $\E_q\sqrt\rho=a$ 得到。
\end{proof}

\section{小 KL 与标准化不能替代矩分析}
\begin{proposition}[KL 趋零而监督方差发散]
\label{prop:small-kl-large-variance}
存在二元正分布族使 $\KL(q\Vert T)\to0$，但
$\Var_q\log(T/q)\to\infty$。也存在 KL 趋零而非负 KL 估计器
$\phi(\rho)=\rho-1-\log\rho$ 方差发散的族。
\end{proposition}
\begin{proof}
第一族取 $q=(\epsilon,1-\epsilon)$，
$T_1=\exp(-\epsilon^{-3/4})$、$T_2=1-T_1$。
第一坐标的 $r=\log(T/q)$ 为
$-\epsilon^{-3/4}-\log\epsilon$，第二坐标为 $O(\epsilon)$。
所以 $-\E_qr\sim\epsilon^{1/4}\to0$，
而二点方差公式给
$\Var_qr=\epsilon(1-\epsilon)(r_1-r_2)^2\sim\epsilon^{-1/2}$。
对固定正 $\kappa$ 和 $\eta<1$，同一族的 $A(\rho_1)$ 趋于有限下界，
$A(\rho_2)=O(\epsilon)$，方差为 $O(\epsilon)$。

第二族仍取 $q_1=\epsilon$，改为 $T_1=\epsilon^{1/4}$。
第一坐标对 KL 的贡献为 $O(\epsilon|\log\epsilon|)$，
第二坐标贡献渐近为 $\epsilon^{1/4}$。
此时 $\rho_1=\epsilon^{-3/4}$、
$\phi(\rho_1)\sim\epsilon^{-3/4}$，
而 $\rho_2=1-O(\epsilon^{1/4})$ 给 $\phi(\rho_2)=O(\epsilon^{1/2})$。
同一二点方差公式再次给 $\Var_q\phi\sim\epsilon^{-1/2}$。
两族说明，均值散度小不能单独控制单动作估计的二阶矩。
\end{proof}

\begin{proposition}[二元总体标准化消去信号形状]
\label{prop:binary-standardization}
设 $q=(v,1-v)$，$0<v<1$，两点信号差 $D=f_1-f_2\ne0$。
令 $\sigma^2=\Var_qf$，取第一 logit 的 score
$\psi=(1-v,-v)$。则
\begin{equation}
\Var_q\!\left[\psi\,\frac{C_qf}{\sigma}\right]=(1-2v)^2.
\end{equation}
任何严格递增的信号变换均保留逐点标准化值。
\end{proposition}
\begin{proof}
有 $C_qf=((1-v)D,-vD)$、
$\sigma=|D|\sqrt{v(1-v)}$。除以 $\sigma$ 后只剩
$\operatorname{sign}(D)$ 和 $v$。两个梯度取值之差为
$\operatorname{sign}(D)(1-2v)/\sqrt{v(1-v)}$，
乘以二点方差系数 $v(1-v)$ 得结论。
严格递增变换保留 $D$ 的符号。本命题使用总体方差、二元动作和单个 logit；
多动作、有限组经验标准化或共享参数不由此式涵盖。
\end{proof}

\section{共同评价指标下的收敛与稳定性}
\label{sec:shape-comparison}
本节固定有限正分布、同一教师及 $0<\eta<1$，比较精确目标映射。相同 $\eta$ 固定插值强度，不保证不同形状移动相同距离；若比较单位 Fisher 速度的瞬时下降，应使用命题~\ref{prop:fisher-shape}。这里的结论也不直接给出共享参数训练的用时排序。

\begin{theorem}[所有形状的正向 KL 乘法收缩]
\label{thm:forward-kl-contraction}
对 $0\le\kappa\le1$，记幂均值的未归一化坐标为 $h$、$Z=\sum_a h_a$。精确目标满足
\begin{equation}
\KL(T\Vert m_{\kappa,\eta})
\le(1-\eta)\KL(T\Vert q)+\log Z
\le(1-\eta)\KL(T\Vert q).
\end{equation}
因而固定教师、固定 $\eta$ 的精确迭代有 $\KL(T\Vert q_n)\le(1-\eta)^n\KL(T\Vert q_0)$。几何端点的第一步不等式取等。
\end{theorem}
\begin{proof}
设 $b=1-\eta$。对 $\kappa>0$，加权算术--几何均值不等式给 $bq_a^\kappa+\eta T_a^\kappa\ge q_a^{b\kappa}T_a^{\eta\kappa}$，所以 $h_a\ge q_a^bT_a^\eta$。幂均值不等式还给 $h_a\le bq_a+\eta T_a$，故 $Z\le1$。几何端点直接使用 $h_a=q_a^bT_a^\eta$，并由加权算术--几何均值不等式得到同一上界。于是
\begin{equation}
\log m_a\ge b\log q_a+\eta\log T_a-\log Z.
\end{equation}
两边从 $\log T_a$ 中减去、乘 $T_a$ 后求和，即得第一式；$\log Z\le0$ 给第二式。逐轮递推完成证明。这个共同的正向 KL 因子，不意味着所有形状对每个凸生成散度都有相同因子；算术端点的全散度乘法收缩是更强结论。
\end{proof}

\begin{lemma}[平方根目标的双锚点恒等式]
\label{lem:root-two-anchor}
固定正教师 $T$，设 $\eta>0$，两个锚点 $q,\widetilde q$ 的平方根未归一化向量均逐坐标为正。令 $b=1-\eta$、$w(q)=b\sqrt q+\eta\sqrt T$、$Z(q)=\|w(q)\|_2^2$，记 $M(q)=m_{1/2,\eta}(q,T)$。则
\begin{equation}
H^2(M(q),M(\widetilde q))
=\frac{b^2H^2(q,\widetilde q)
-[\sqrt{Z(q)}-\sqrt{Z(\widetilde q)}]^2}
{\sqrt{Z(q)Z(\widetilde q)}}.
\label{eq:root-two-anchor}
\end{equation}
\end{lemma}
\begin{proof}
正性保证 $\sqrt{M(q)}=w(q)/\|w(q)\|$。对两个非零向量 $w,x$，展开内积得到
\begin{equation}
\|w-x\|^2=(\|w\|-\|x\|)^2
+\|w\|\|x\|\left\|\frac{w}{\|w\|}-\frac{x}{\|x\|}\right\|^2.
\end{equation}
代入 $x=w(\widetilde q)$，使用 $w(q)-w(\widetilde q)=b(\sqrt q-\sqrt{\widetilde q})$，移项即得恒等式。平方根正性是该解释的条件，不能用负坐标平方后得到的另一映射代替。
\end{proof}

\begin{proposition}[插值区间的扰动收缩比较]
\label{prop:interpolation-stability}
固定 $0<\eta<1$、$b=1-\eta$。三种形状满足
\begin{align}
d_H(M_0(q),M_0(\widetilde q))&=b\,d_H(q,\widetilde q),\\
H(M_{1/2}(q),M_{1/2}(\widetilde q))&\le\frac b{\sqrt{b^2+\eta^2}}H(q,\widetilde q),\\
\TV(M_1(q),M_1(\widetilde q))&=b\,\TV(q,\widetilde q),\\
H(M_1(q),M_1(\widetilde q))&\le\sqrt b\,H(q,\widetilde q).
\end{align}
平方根的上述 Hellinger 上界常数不大于算术上界常数，当且仅当 $\eta\ge1/2$。
\end{proposition}
\begin{proof}
几何目标的 log-ratio 差为 $b\log(q/\widetilde q)$ 加公共常数，取跨度得到第一式。平方根目标的 $Z(q)=b^2+\eta^2+2b\eta\sum_a\sqrt{q_aT_a}\ge b^2+\eta^2$；在式~\eqref{eq:root-two-anchor} 中丢弃被减去的平方项并开方得到第二式。算术目标差精确等于 $b(q-\widetilde q)$，故 TV 等式成立。

对每个坐标，Cauchy--Schwarz 给
\begin{equation}
\sqrt{(bq_a+\eta T_a)(b\widetilde q_a+\eta T_a)}
\ge b\sqrt{q_a\widetilde q_a}+\eta T_a.
\end{equation}
求和并用 $H^2=2-2\sum\sqrt{pq}$ 得最后一式。比较两正上界常数并平方，$b^2/(b^2+\eta^2)\le b$ 等价于 $b\le b^2+\eta^2$，即 $\eta\ge1/2$。这里只比较已证明的统一界，不推断任意具体锚点对有相同排序。
\end{proof}

\begin{proposition}[平方根插值的角度递推]
\label{prop:root-interpolation-angle}
固定 $0<\eta<1$，令 $\theta=\arccos\sum_a\sqrt{q_aT_a}\in[0,\pi/2)$。目标的夹角满足
\begin{equation}
\theta^+=\arctan\frac{(1-\eta)\sin\theta}{\eta+(1-\eta)\cos\theta}.
\end{equation}
在 $\eta=1/2$ 时，$\theta^+=\theta/2$，因而 $H^2(q_n,T)=4\sin^2(\theta_0/2^{n+1})$。
\end{proposition}
\begin{proof}
当 $q\ne T$，写 $\sqrt q=\sqrt T\cos\theta+w\sin\theta$，其中 $w$ 为与 $\sqrt T$ 正交的单位向量。未归一化目标的两个正交系数分别为 $\eta+(1-\eta)\cos\theta$ 和 $(1-\eta)\sin\theta$。前者为正，二者之比即夹角正切。$\eta=1/2$ 时用 $\sin\theta/(1+\cos\theta)=\tan(\theta/2)$ 得精确二分；迭代后使用 $H^2=4\sin^2(\theta/2)$。$q=T$ 时连续延拓给零角度。
\end{proof}

\section{形状优劣依赖比较准则}
\label{sec:shape-criteria-summary}
三个端点分别在 log-ratio、平方根和概率坐标中插值。它们共享定理~\ref{thm:forward-kl-contraction} 的正向 KL 全局界与局部 $(1-\eta)^2$ 的 KL 比例，但只在算术端点直接得到所有凸生成散度的同因子收缩。教师覆盖系数分别为：几何端点无统一正下界，平方根至少 $\eta^2T_a$，算术至少 $\eta T_a$。

监督信号采用相同标准化 $A=f/\eta$。其小步极限依次为 $\log\rho$、$2(\sqrt\rho-1)$、$\rho-1$。平方根方差为 $4(1-c^2)\le4$，其中 $c=\sum\sqrt{qT}$；算术方差为 $\chi^2(T\Vert q)$，几何对数方差也无统一有限上界。固定正 $\eta<1$ 的算术信号仍有有限统一界，不能与其最坏方差在 $\eta\downarrow0$ 时发散混淆。这些矩结论及最优形状的不同准则已经在本附录证明；它们约束采样信号，不直接给出实际优化时间排序。

\begin{proposition}[共同 KL 指标不存在固定端点排序]
\label{prop:shape-order-counterexample}
在 $\eta=1/2$ 下，存在正二动作分布使三种形状的更新后逆向 KL 按 $1,1/2,0$ 从小到大排列，也存在按 $0,1/2,1$ 排列的分布。因此平方根不能在这个共同指标下逐实例支配另外两个端点。
\end{proposition}
\begin{proof}
给定 $q=(x,1-x)$、$T=(t,1-t)$，记目标第一坐标为 $p_\kappa$。三式分别为
\begin{align}
p_0&=\frac{\sqrt{xt}}{\sqrt{xt}+\sqrt{(1-x)(1-t)}},\\
p_{1/2}&=\frac{(\sqrt x+\sqrt t)^2}
{2+2[\sqrt{xt}+\sqrt{(1-x)(1-t)}]},\qquad p_1=\frac{x+t}{2}.
\end{align}
第一组 $x=1/20,t=1/5$ 代入后有 $0.10290<p_0<0.10291$、$0.11411<p_{1/2}<0.11412$、$p_1=0.125<t$。第二组 $x=3/5,t=19/20$ 有 $p_1=0.775$、$0.80678<p_{1/2}<0.80679$、$0.84223<p_0<0.84224<t$；这些有理数夹界由对正根式两边平方即可核对。二元 KL 对 $p$ 的导数是 $\log[p(1-t)/(t(1-p))]$，在 $p<t$ 时严格为负，故两组概率次序给出相反的严格 KL 排序。代入二元 KL 后的近似值如下。
\begin{center}
\begin{tabular}{ccrrr}
\toprule
$x$ & $t$ & $\kappa=0$ & $\kappa=1/2$ & $\kappa=1$\\
\midrule
$0.05$ & $0.2$ & $0.034382$ & $0.026311$ & $0.019660$\\
$0.6$ & $0.95$ & $0.079878$ & $0.129358$ & $0.180628$\\
\bottomrule
\end{tabular}
\end{center}
两个严格反向排序否定固定排序；相同 $\eta$ 并没有使各形状移动相同距离，故该例不替代单位 Fisher 速度的命题~\ref{prop:fisher-shape}。
\end{proof}

在可实现教师附近，有限 BPO、MDPO 的共同参数谱由定理~\ref{thm:joint-spectrum} 描述。平方根的采样矩控制不能改写为改善了固定参数 Fisher 矩阵的最小特征值。选择形状时须分别说明评价散度、动作积分方式、内层预算和优化器；这些条件不同，最优性判断也不同。

